\section{Exact degrees and highest homogeneous forms}
\label{sec:saturated-two-root}

This section proves the global degree bounds and their attainment for
the four literal circuits.  Write $\gamma=\varepsilon$ for the fixed
center choice; it is not another unknown.  The square-recovery argument
and the complete witness constructions were proved in the preceding
section.

\subsection{Literal degrees and highest homogeneous forms}

Set $q_0=\beta J$ and $k_0=\eta+\zeta$.  With the port values fixed
as constants, the four exact results are
\[
\begin{array}{c|r|r}
\text{circuit}&\text{live binary operations}&\text{total degree}\\\hline
 X=wq,\ C_1=(A+1)UE&92&283\\
 X=wq,\ C_0=AUE&93&278\\
 X=q+x,\ C_1=(A+1)UE&98&258\\
 X=q+x,\ C_0=AUE&99&254.
\end{array}
\]
The operation counts refer to the supplied literal graphs after
discarding unused nodes; they charge additions, subtractions and
multiplications once each, including multiplication by a fixed integer.

\paragraph{Exact-degree certificate.}
For the $X=wq$ circuits, the highest homogeneous forms are
\begin{align*}
 [N_{2,1}]_{283}
   &=2^{11}q_0^{149}w^{29}Y^{61}i^{10}k_0^{32}(q_0-F)^2,\\
 [N_{2,0}]_{278}
   &=2^{10}q_0^{146}w^{28}Y^{60}i^{10}k_0^{32}(q_0-F)^2.
\end{align*}
For the direct circuits, put $x_0=q_0+x$.  The highest forms are
\begin{align*}
 [N_{2,1}^{\rm dir}]_{258}
   &=2^{11}q_0^{124}x_0^{29}Y^{61}i^{10}k_0^{32}(q_0-F)^2,\\
 [N_{2,0}^{\rm dir}]_{254}
   &=2^{10}q_0^{122}x_0^{28}Y^{60}i^{10}k_0^{32}(q_0-F)^2.
\end{align*}
These are nonzero polynomials for every fixed $\beta>0$.  Therefore
they certify exact degrees for the actual coefficient tuple.

\paragraph{Proof.}
Scale all ten witnesses by $t$, keep the coefficient ports fixed, and
work in the formal Laurent-series field at $t=\infty$ over the rational
function field in the ten unscaled witnesses.  A series degree is the
largest exponent of $t$.  For polynomials it is their ordinary total
degree.  Choosing formal square roots with the indicated square leading
forms is justified by the binomial series; each relative correction
has negative degree.  No analytic convergence is used.

In either source the literal numerator satisfies
\[
 Z:=Z_{\rm num}/M=V+C_\gamma+\mathcal E,
\]
where the rational remainder has degree at most $4$.
For $X=wq$, it is
\[
 \mathcal E=n-LB_t+K_{\rm pos}/H+Y/q^3-B_t/(q-1)-e/\mu^2,
 \quad L=k(T+\alpha),\quad H=4A-5,
\]
and $M=q^3(q-1)\mu^2H$.
For $X=q+x$, replace the remainder by
\[
 \mathcal E=n-LB+K_{\rm pos}/H+Y/q^3
                -B/[q(q-1)]-X/q^4-e/\mu^2,
\]
and retain the extra factor of $q$ in
$M=q^4(q-1)\mu^2H$.  Both denominators have degree $13$.
These rational identities are exact readings of the respective
polynomial graphs, not additional defining equations.

Choose formal roots $g^2=ED_f$ near $AUE$ and
$y^2=D_y/\delta$ near $V$.  The two exact identities are
\[
 (AUE)^2-ED_f=U^2E^2-E,                         \tag{D1}
\]
and
\[
 (y-V)(y+V)=(V^2-1)/\delta.                       \tag{D2}
\]
The degrees controlling the argument are
\[
\begin{array}{c|rrrrrrrrrr}
 &A&U&E&D_f&\delta&V&M&g&C_1-g&C_0-g\\\hline
 X=wq&5&5&18&38&48&42&13&28&23&18\\
 X=q+x&4&5&16&34&42&38&13&25&21&17.
\end{array}
\]
Equation (D2) gives $\deg(y-V)=-6$ and $-4$, respectively.
Equation (D1) gives
\[
 C_0-g=\frac{U^2E^2-E}{AUE+g},
 \qquad C_1-g=UE+(C_0-g).
\]
The leading forms of these two defects are $u_0e_0/(2a_0)$ and
$u_0e_0$, respectively, where the subscripts denote leading forms.
In every case the defect degree is greater than $4$, so it dominates
both $\mathcal E$ and $V-y$.

The exact four-factor norm identity is
\[
 N_2=\delta^2M^4
       \prod_{\sigma,\rho\in\{-1,1\}}(Z+\sigma y+\rho g).
                                                               \tag{D3}
\]
The two factors with $\sigma=1$ have leading form $2v_0$.
The factor with $\sigma=-1,\rho=1$ has leading form $2g_0$.
The last factor has leading form $u_0e_0$ for $C_1$ and
$u_0e_0/(2a_0)$ for $C_0$.
Thus the four degree totals are exactly
\begin{align*}
 2(48)+4(13)+2(42)+28+23&=283,\\
 2(48)+4(13)+2(42)+28+18&=278,\\
 2(42)+4(13)+2(38)+25+21&=258,\\
 2(42)+4(13)+2(38)+25+17&=254.
\end{align*}

To verify the coefficients without any large-numeral expansion, for
$X=wq$ put
\[
 a_0=q_0^3wY,\quad u_0=ik_0^2Y^2,\quad
 e_0=q_0^{10}w^2Y^4k_0^2,\quad r_0=q_0^3(q_0-F).
\]
Then
\[
 d_0=a_0^4u_0^2e_0,\quad m_0=4q_0^8a_0,\quad
 v_0=-a_0^2u_0^2e_0r_0,\quad g_0=a_0u_0e_0.
\]
Multiplying the leading factors in (D3) gives
\[
 8d_0^2m_0^4v_0^2g_0u_0e_0
 \quad\hbox{or}\quad
 4d_0^2m_0^4v_0^2u_0^2e_0^2,
\]
respectively.  These simplify to the two asserted $X=wq$ forms.
For the direct source replace
\[
 a_0=q_0^2x_0Y,\qquad
 e_0=q_0^8x_0^2Y^4k_0^2,\qquad m_0=4q_0^9a_0;
\]
the remaining abbreviations retain their formulas.  The same products
give the two direct forms.  Every displayed final top form is a nonzero
polynomial, proving attainment as well as the upper bound.
\hfill$\square$

\paragraph{Independent exact polynomial replay.}
The standard-library verifier evaluates every gate of each supplied
graph as a univariate polynomial over the integers.  On the ray where
all ten witnesses equal $t$ and
$(\beta,K_0,M_C,M_F)=(2,3,5,7)$, it obtains, in table order,
\[
 2^{192}t^{283},\qquad 2^{188}t^{278},\qquad
 2^{167}3^{29}t^{258},\qquad 2^{164}3^{28}t^{254}
\]
as the exact leading terms.  This is finite arithmetic verification of
one specialization.  The global Laurent-series proof, including its
explicit nonzero top homogeneous forms, is what establishes the claims
for the real fixed coefficient tuple.
